\section{Event Selection}
\label{sec:selection_effect}

Event selection is a central part of any hadron-collider analysis. The raw candidate sample contains true signal decays, partially reconstructed backgrounds, misidentified backgrounds, and random combinations of charged tracks. The goal of the baseline selection is to retain a useful signal-rich region while keeping the cuts simple enough to interpret.

The open-data ntuple is already preselected: it contains only candidates with
$\chi^2_{\mathrm{vertex}}(B) < 12$ and $d_{\mathrm{flight}}(B) > 3~\mathrm{mm}$.
The baseline selection applied on top of that is
\begin{align}
  \chi^2_{\mathrm{vertex}}(B) &< 5, \\
  d_{\mathrm{flight}}(B) &> 20~\mathrm{mm}, \\
  P_K(H_1), P_K(H_2), P_K(H_3) &> 0.5 .
\end{align}
Here $\chi^2_{\mathrm{vertex}}(B)$ is the $\chi^2$ of the three-track vertex fit,
not divided by its number of degrees of freedom; $d_{\mathrm{flight}}(B)$ is the
reconstructed distance in millimetres between the primary and secondary vertices;
and $P_K(H_i)$ is the kaon PID response for daughter track $H_i$, a
probability-like variable between 0 and 1.

The selection is severe: it keeps $1645$ of the $5\,135\,823$ MagnetDown
candidates and $1294$ of the $3\,420\,295$ MagnetUp ones, a retention of
$3$--$4\times10^{-4}$. That is a deliberate trade of efficiency
for purity, and it has a consequence used later in this report: after these cuts
there is little residual background left for a classifier to separate.

The vertex-quality cut suppresses candidates whose three tracks do not form a well-reconstructed secondary vertex. This is physically motivated because the daughters of a true B decay should originate from a common decay point. A poor vertex fit is more likely to arise from accidental combinations of tracks produced in different processes or from poorly reconstructed tracks.

The flight-distance selection exploits the finite lifetime of B mesons. Because B hadrons travel a measurable distance before decaying, their decay vertices are displaced from the primary proton--proton interaction point. Requiring a sufficiently large flight distance therefore suppresses prompt background combinations that originate near the primary vertex.

The PID selection implements the kaon mass hypothesis used in the invariant-mass reconstruction. Requiring all three tracks to have kaon-compatible PID values improves the consistency of the reconstructed final state. In a more complete analysis, one would optimize PID requirements with calibration samples and evaluate the corresponding efficiency and misidentification rates. Here the selection is chosen as a transparent demonstration of how detector-level PID information enters the analysis.

The selected sample is used consistently across the mass reconstruction, signal/background classification, CP-asymmetry diagnostic, and Dalitz-plot studies. This consistency is important: changing the selection between notebooks would make comparisons ambiguous and could introduce artificial differences in the resulting distributions.
